\section{Comparison with Schubert's preprint}\label{sec:changes}
We list how the statements above relate to those of \cite{S}, equations (1.15)--(1.36) there.
\begin{itemize}
\item \emph{The gap.} In the proofs of \cite[Theorems~4.1 and~4.4]{S}, the steps (4.16)$\to$(4.17), (4.55)$\to$(4.56) and
(4.62)$\to$(4.63) infer $\norm{\ue^{A}a}\le\norm{\ue^{B}a}$ from an operator inequality $A\le B$ between non-commuting
self-adjoint operators. The exponential is not operator monotone: with $A=\mathrm{diag}(0,1)$,
$P=\begin{pmatrix}1&1\\1&1\end{pmatrix}\ge0$ and $v=(-3,1)$, the function $g(s)=\norm{\ue^{A+sP}v}^2$ has
$g'(0)=4(6-\ue^2)<0$. In addition, \cite[Def.~1.1]{S} and the proofs of \cite[Prop.~3.5, Lemma~3.11]{S} take the two weights
of the norm in different orders. The proof of \cite[Lemma~3.11]{S} passes from one order to the other by
\cite[Proposition~7.1, (7.2)]{S}, which states that the two norms are equivalent; the sketch of its proof rests on
\cite[Theorem~7.2]{S}, which is stated for small parameters $|\alpha|,|\beta|\le c$. For $\beta>\tfrac12$, the range the theorems
need, the two norms are not equivalent (Proposition~\ref{pc:classes}).
\item \emph{Amplitude class.} Theorem~\ref{pc:main} uses the norm of \cite[Def.~1.1]{S} for every bounded range $0<\alpha\le A$, with constants depending on $A$.
Theorem~\ref{thm:dd14} and its corollaries treat the reversed order used in Schubert's proofs.
\item \emph{\cite[Theorem~1.3]{S}.} Its estimate \cite[(1.18)--(1.19)]{S} states the error $C\hbar\norm a_{\alpha,\beta}$ for $0\le t\le\delta\hbar^{-1/2}$;
Section~\ref{subsec:corrections} shows that the error has the order $\hbar t$. The introduction of \cite{S} describes it in the same
way: the dispersion ``takes place on a scale of order $\hbar t$, and this is responsible for the error term in (1.18)'', and the
expression \cite[(1.26)]{S} ``is of order $\hbar t$''. The corrected statement is Remark~\ref{thm:main}, with the same window and
generality in $\Gamma$, for every $\beta>\tfrac12$ instead of $\beta>1$, and also for the reversed-weight norm.
\item \emph{Sobolev bounds on the quotient.} The inequalities \cite[(3.3), (3.4)]{S} bound $u_\Gamma$ in $L^2(M)$, respectively pointwise, by the $L^1(M)$
norms of $u_\Gamma$ and $\Delta^su_\Gamma$, for every $\Gamma$, and \cite[Proposition~3.2 and Corollary~3.4]{S} combine them with the bound
$\norm{u_\Gamma}_{L^1(M)}\le\norm u_{L^1}$. On a quotient with a cusp both Sobolev bounds fail (Remark~\ref{rem:g4}). Here periodization is
controlled in weighted $L^2$ spaces (Lemma~\ref{lem:t2}, Proposition~\ref{os:p1}), and pointwise bounds use the local Sobolev inequality
with the factor $m(z)^{1/2}$ (Lemma~\ref{pc:g43}).
\item \emph{\cite[Theorem~1.4]{S}.} Its estimate \cite[(1.30)--(1.31)]{S} states the bound
\[C^{N+K+1}N^{2N}K!\,(\alpha\hbar)^{-2N-4}(\hbar t/\alpha^2)^K\norm a_{\alpha,\beta}\]
for $\alpha^3\ge\beta^2\hbar$ and $0\le t\le\delta\alpha^{3/2}\hbar^{-1/2}$. Theorem~\ref{pc:main} retains the dependence $C^{N+K+1}N^{2N}K!$ on the orders, with the condition
$\alpha^3\ge\kappa\hbar$ for any fixed $\kappa>0$ ($\kappa=\beta^2$ in \cite{S}), the longer window
$t\le\delta(\alpha/\hbar)^{1/2}$, and a bound that is smaller than that of \cite{S} by the factor $(\alpha\hbar)^4\hbar t/\alpha^2$,
which is at most $1$ on the window of \cite{S}.
Theorem~\ref{pc:main} covers every bounded range of radii, with constants depending on its upper bound.
\item \emph{\cite[Corollary~1.5]{S}.} The range of $K$ stated in \cite[(1.32)]{S} is replaced by the multiplicative range
$K\rho\in[1-\theta,1-\theta/2]$ of Corollary~\ref{pc:cor15}, which is the one used in the proof in \cite{S}. The prefactor is $(\alpha\hbar)^{-2N}$
instead of $(\alpha\hbar)^{-2N-4}$.
\item \emph{\cite[Corollary~1.6]{S}.} Its estimate \cite[(1.34)--(1.36)]{S} states the pointwise bound
\[C^{N+K+1}(2N)!\,K!\,(\alpha\hbar)^{-2N-8}(\hbar t)^K\norm a_{\alpha,\beta}\]
for every $\Gamma$. The specialization of Corollary~\ref{pc:cor16} to a positive
injectivity radius has the prefactor $(\alpha\hbar)^{-2N-2}$ and the time factor $(\hbar t/\alpha^2)^{K+1}$. For fixed
$\alpha$ the factor $\alpha^{-2K}$ can be put into $C^K$, and the time factor of \cite{S} follows; with constants uniform in
$\alpha$ the factor $\hbar t/\alpha^2$ is what Theorem~\ref{pc:main} gives. On surfaces with cusps the bound holds with the factor
$m(z)^{1/2}$ (Corollary~\ref{pc:cor16}). The prefactor $(\alpha\hbar)^{-2N-8}$ of \cite{S} absorbs this factor where
$m(z)\le(\alpha\hbar)^{-12}$, since there $m(z)^{1/2}(\alpha\hbar)^{-2N-2}\le(\alpha\hbar)^{-2N-8}$; higher in the cusp the factor cannot be
removed on the modular surface (Theorem~\ref{cl:necessary}, Remark~\ref{cl:reading}). In \cite[(1.36)]{S}, $\psi^{(K)}$ should read $u^{(K)}$.
\end{itemize}
